\subsection{Data Collection}

We collected all of our data from a sandbox that we wrote in Python. The sandbox stores the five redesigned laws as hidden functions: the spring $F = -5\,|x|^{1.5}\,\mathrm{sign}(x)$, kinetic energy $E_k = 2/v$, wave speed $v = 15\ln(0.5f)$, the screened Coulomb force $F = 6e^{-0.4r}/r^2$ and the saturating Ohm's law $I = 5\tanh(V/2)$. A model never sees these formulas. It is only told what the system is, which constants are fixed (for example a cart mass of 2\,kg), what range of inputs it may use, and the readings that come back.

What a model is allowed to observe depends on the sandbox design, and we used four of them.
\begin{itemize}
  \item \textbf{Design S} gives the model one fixed table of readings in a single prompt. We generated 50 tables for each law so that every model answers on the same data. A spring table is a simulated release of a 1\,kg block over 10\,s with a 0.1\,s time step. Tables for the other laws hold 50 readings spread over the input range.
  \item \textbf{Design A} lets the model pick an input and read the exact output, up to 12 times.
  \item \textbf{Design B} is only for the spring. The model releases the block and sees how it moves, up to 6 times, and has to work out the force from the motion.
  \item \textbf{Design C} works like design A with up to 10 measurements, but each reading has Gaussian noise with a standard deviation of 3\% of its value.
\end{itemize}

In designs A, B and C the model works alone. It gets one objective prompt, sends one action at a time, reads the result and finally submits a formula. We did not step in during any run. Each run has a seed that fixes the noise, so we can repeat it exactly.

We tested nine models through the RapidScreen API: gpt-6-sol, gpt-6-luna, gpt-5.6-terra, gpt-5.5, gpt-5.4, gpt-5.4-mini, gpt-5.4-nano, deepseek-v4-flash and deepseek-v4.1-flash. Early in the project we also ran two small local models with Ollama (qwen2.5-coder:7b and qwen2.5:3b) on the spring in design S. A power law curve fit and a plain Hooke's law guess were used as baselines that involve no language model.

Table~\ref{tab:datacount} shows how much data we ended up with. Every model did 10 runs per law in designs A and C, and 10 runs on the spring in design B. For design S, all nine models answered the 50 kinetic energy tables. Our API key was replaced after that and five of the models were not included in the new one, so the rest of design S only has the four models that were on both keys.

\begin{table}[h]
\centering
\small
\begin{tabular}{lll}
\hline
Data & Coverage & Runs \\
\hline
Design A & 9 models $\times$ 5 laws $\times$ 10 seeds & 450 \\
Design C & 9 models $\times$ 5 laws $\times$ 10 seeds & 450 \\
Design B & 9 models $\times$ spring $\times$ 10 seeds & 90 \\
Design S & 9 models on kinetic energy, 4 models on the other laws, 50 tables each & 1{,}250 \\
Self improvement & gpt-5.4-nano against gpt-5.5, tested on each law & 5 laws \\
\hline
\end{tabular}
\caption{Data collected for the benchmark.}
\label{tab:datacount}
\end{table}

For the self improvement test, gpt-5.4-nano first practised in design C on four training springs that are not among the five test laws. It rewrote a short set of strategy notes after each of three rounds. We then ran it 10 times on the test law without the notes and 10 times with them, and compared both against 10 runs of gpt-5.5.

Each run is saved as a JSON log that holds the prompt, every action and reading, the final answer, the grade, the time taken and the token count. Some runs failed because the internet or the API connection dropped, for instance when the laptop went to sleep. Those runs never reached the model, so we deleted them and ran them again. Runs that the provider's content filter blocked were kept as failures. Where a model had more than 10 runs for a law and design, we used the 10 latest.

Our data has some weaknesses. It is simulated, so it has none of the drift, friction or sensor faults of a real lab. All five laws have one input, and results could be different for laws with several variables. Every model got the same prompt wording, and other wording might shift some scores. Ten runs can show a large gap between two models, but a gap of one or two runs should not be treated as meaningful.
